\section{Exact correspondence with the Lean formalization}
\label{sec:formal}

The statements of Section~\ref{sec:model} are the paper form of named Lean~4 declarations of the
public development~\cite{LeanRepo}; the model (the rotations, the evolution and
\path{Admissible}), the cost $T=\texttt{cost}\,\theta$ and the carrier
$h$ of~\eqref{eq:h} are fixed in \path{RobustZ/Statement.lean}, and the scalar flatness~\eqref{eq:hflat} is derived from the matrix flatness rather than assumed. Write $q=2N+2$,
$\tauvar=T/2$ and $k=2N+1=q-1$; the
sharpened layer uses $\rho_a(q)=1-a(\log q/q)^{2/3}$, the exponent $\delta'>0$ with
$\cosh\delta'=1/\rho_a(q)$, the band approximation error
$\eta=\texttt{epsOf}\ \tauvar\,\delta'\,k$, and the collar half-width $B_\sharp=\sqrt\eta$. The endgame
statements are stated for $\phi=\pi$, the case formalized.

\subsection{The machine-checked statements}

\begin{theorem}[Elementary bound, all orders; \path{elementary_bound}, \path{elementary_bound_pi}]
\label{thm:formal-elementary}
For every $0<\phi\le\pi$ and every order-$N$ admissible sequence,
\begin{equation}
T\ge2\,\bigl(2\,q!\,\sin^{2}(\phi/4)\bigr)^{1/q},
\qquad\text{and at }\phi=\pi:\quad T\ge2\,(q!)^{1/q}.
\label{eq:formal-elementary}
\end{equation}
Moreover
(\path{elementary_gt_four_over_e}) $4(N+1)/\e<2\,(q!)^{1/q}$ for every $N\ge0$, so the
bound implies, and is strictly stronger than, the closed form $4(N+1)/\e$ of
Section~\ref{sec:elementary}.
\end{theorem}

\begin{theorem}[Sharpened bound; \path{endgame_bound}]
\label{thm:formal-sharp}
For every $a\ge2$ and every $N\ge N_0(a)$, with $N_0(a)$ the explicit integer threshold~\eqref{eq:formal-bookkeeping}, every order-$N$ admissible sequence with $\phi=\pi$ satisfies
\begin{equation}
4\,(N+1)\Bigl(1-a\Bigl(\frac{\log q}{q}\Bigr)^{2/3}\Bigr)\le T,
\label{eq:formal-sharp}
\end{equation}
and the two numerical instances \path{endgame_bound_2_3_num} and
\path{endgame_bound_2_2_num} bound it explicitly:
\begin{equation}
\begin{aligned}
N\ge5010\ &\Longrightarrow\
4(N+1)\bigl(1-2.3\,(\log q/q)^{2/3}\bigr)\le T,\\
N\ge46000\ &\Longrightarrow\
4(N+1)\bigl(1-2.2\,(\log q/q)^{2/3}\bigr)\le T .
\end{aligned}
\label{eq:formal-sharp-num}
\end{equation}
The formal inequality is~\eqref{eq:formal-sharp}, and it is non-strict: the proof rules out the non-strict case
$\tauvar\le\rho_a(q)\,q$, so strictness is not part of the kernel-checked statement. Combined with
the elementary term (\path{all_N_bound_end}), the same layer gives, for every $N\ge1$,
\begin{equation}
\max\Bigl(2\,(q!)^{1/q},\ \mathbf 1_{N\ge N_0(a)}\cdot
4\,(N+1)\bigl(1-a(\log q/q)^{2/3}\bigr)\Bigr)\le T,
\label{eq:formal-allN}
\end{equation}
which is the kernel-checked form of~\eqref{eq:allN}.
\end{theorem}

The threshold is the explicit function \path{effN0End}, evaluated at the cutoff constant
$K_0=21.1$ and at $(C_{d1},C_{d2},\beta)=(7,770,4)$:
\begin{equation}
\begin{aligned}
&N_0(a)=\max\bigl(\lceil256a^{3}\rceil,\ \lceil\texttt{effQ0}/2\rceil,\
\lceil\texttt{effQmin}/2\rceil\bigr),\quad
\texttt{effQ0}=\max\bigl(512a^{3},(64\texttt{effK1})^{-1/\texttt{effE}}\bigr),\\
&\texttt{effE}=\max\bigl(\beta+\tfrac23-2s,\ \tfrac32-\tfrac32s\bigr),\qquad
s=\texttt{effS}(a)=\tfrac{47}{50}a^{3/2},\\
&\texttt{effK1}=\frac{1+C_{d1}}{36}\Bigl(2704\,C_{d2}/a+5408\,K_0C_{d1}/a+752\,K_0\,a^{-3/4}\Bigr),\\
&\texttt{effQmin}=\max\bigl(10020,\ \e,\ (20\,\e\,a^{-1/4})^{1/(s/2-7/6)}\bigr).
\end{aligned}
\label{eq:formal-bookkeeping}
\end{equation}
The general threshold \path{effN0} of \path{Effective.lean} carries the crude floor
$5\cdot10^{9}$, which \path{effN0End} replaces by \path{effQmin}. The numerical lemmas behind~\eqref{eq:formal-sharp-num} give
$-\texttt{effE}(2.3,4)\ge1.85$, $\texttt{effK1}(2.3)\le280249$,
$\texttt{effQ0}(2.3)\le10020$, $\texttt{effQmin}(2.3)\le10020$ and hence $N_0(2.3)\le5010$; and
$-\texttt{effE}(2.2,4)\ge22/15$, $\texttt{effK1}(2.2)\le292954$,
$\texttt{effQ0}(2.2)\le92000$, $\texttt{effQmin}(2.2)\le92000$ and hence $N_0(2.2)\le46000$ (the
source records the sharper value $\texttt{effN0End}\ K_0\,2.2=45006$ and rounds up).

\subsection{The constants}

The real-analytic inputs are proved in \path{RobustZ/EffectiveCalc.lean}: for the contour-shift
exponent $\mu(\rho)=\log\bigl((1+\sqrt{1-\rho^{2}})/\rho\bigr)-\sqrt{1-\rho^{2}}$ and its prefactor
$D_0$, the sharpened layer uses $\mu(\rho)\ge\frac{47}{50}(1-\rho)^{3/2}$ on $\frac12\le\rho<1$ and
$D_0(\rho)\le13\,(1-\rho)^{-1/2}$ on $\frac34\le\rho<1$ (the asymptotic constant $2\sqrt2/3$
leaves $47/50$ within $0.3\%$ of the optimum). The criterion
$D_0(\varrho_q)\e^{-q\mu(\varrho_q)}<\sqrt{\varrho_q}/(Cq)$ packages the two at threshold
$q_0(C)=\max(10^{4},7\max(C,100)^{6/7})$, $C>0$; the endgame path uses instead the combined form
\path{D0_mu_le}, $D_0(\rho_a(q))\e^{-q\mu(\rho_a(q))}\le13\,a^{-1/2}(q/\log q)^{1/3}q^{-s}$,
with threshold \path{effQ0}.

\begin{table}[t]\centering\small
\caption{Explicit constants of the machine-checked sharpened bound and the declarations carrying
them; the last two rows are the collar constants used by the final assembly.}
\label{tab:formal-constants}
\begin{tabular}{@{}>{\raggedright\arraybackslash}p{0.62\linewidth}>{\raggedright\arraybackslash}p{0.32\linewidth}@{}}
\toprule
\textbf{Constant} & \textbf{Lean declaration}\\
\midrule
$\mu(\rho)\ge\frac{47}{50}(1-\rho)^{3/2}$ & \path{mu_lower_sharp}\\
$D_0(\rho)\le13\,(1-\rho)^{-1/2}$ & \path{D0_le_inv_sqrt}\\
$q_0(C)=\max(10^{4},7\max(C,100)^{6/7})$ & \path{criterion_param_max}\\
$K_0=21.1$; $|\chi'|\le K_0/B$, $|\chi''|\le K_0/B^{2}$ & \path{deriv_cutoff_le_explicit}\\
$\norm{\Psi'}\le(7k^{2}/\tauvar+9)\eta$; $\norm{\Psi''}\le(384k^{4}/\tauvar^{2}+14k^{2}/\tauvar+9)\eta$ & \path{norm_deriv(2)..._sharp}\\
$\tauvar M_1\le7k^{2}\eta$, $\tauvar^{2}M_2\le770k^{4}\eta$ & \path{collar_bound_concrete}\\
\bottomrule
\end{tabular}
\end{table}
The abbreviations \path{norm_deriv(2)..._sharp} and \path{..._collar_gen} denote the
\path{sharp} and \path{gen} instances of
\path{norm_deriv_psiA_scaled_}\hspace{0pt}\path{collar_*} and
\path{norm_deriv2_psiA_scaled_}\hspace{0pt}\path{collar_*}.

The sharp collar bounds hold on the whole collar $|\mu|\le\tauvar+B$ under explicit smallness
conditions on $B/\tauvar$ and on the tail ratio $r$. The generalized forms used by the final
assembly (\path{..._collar_gen}) keep the honest factor $R=1/(1-r')$ instead of assuming it
small, where $\eta_h$ is the honest constant of the coefficient bound
$|c_{k+j}(\tauvar)|\le\eta_h\,r^{j}$ (\path{fcoef_tail_le}) and
$r'=r\,\e^{2\sqrt{B/\tauvar}}$:
$\norm{\Psi'}\le2\,\e^{2k\sqrt{B/\tauvar}}(2k^{2}R+4R^{3}+\tauvar R)\,\eta_h/\tauvar$, with a
five-term analogue for $\Psi''$ (the Lean function \path{psiA}); the exact cancellation
\begin{equation}
\eta_h=\e^{\tauvar\sinh\delta'}\e^{-\delta'k}\le\frac{\eta}{2R}
\qquad(\texttt{honest\_le\_eps\_mul})
\label{eq:formal-cancel}
\end{equation}
then removes $R$ and leaves the constants $7$ and $770$, free of $q$ and of $\tauvar$. The factor
$\frac1{36}(1+C_{d1})$ of~\eqref{eq:formal-product} is the elementary estimate
$(2\pi)^{-2}<1/36$; at the end of the chain the three terms of~\eqref{eq:formal-product} are each
bounded by $\texttt{effK1}\,q^{\texttt{effE}}(\log q)^{-1/2}$, and $q\ge\texttt{effQ0}$ makes their
sum smaller than $1/64$. This is where $a\ge2$ enters: $s=\frac{47}{50}a^{3/2}$ must exceed
$\beta/2+1/3$, that is $a>1.834$ for $\beta=4$; $a=2$ is a safe margin, not a sharp one.

\subsection{Mathematical content of the formalization}

Three ingredients are worth naming. First, the doubling of flatness: order-$N$ matrix robustness at $\lambda=1$ forces $h^{(i)}(1)=0$ for
$i<q=2N+2$ (\path{h_flat}), so the $q$ vanishing derivatives become $q$ vanishing moments
(\path{h_moments}), and \path{abs_zero_le_of_flat} turns Taylor--Lagrange with $q$ flat
terms into $|h(0)|\le\tauvar^{q}/q!$. Second, the collar estimates cover the whole collar
$|\mu|\le\tauvar+B_\sharp$ with $B_\sharp=\sqrt\eta$, beyond the band where the polynomial only has to
approximate $\e^{-i\mu}$; the version with honest $1/(1-r')$ constants is what makes a width of
order $\sqrt\eta$ usable with no unproved hypothesis. Third, the product bound is sharpened from
the $q^{8}$ estimate of Section~\ref{sec:wholeband} to
\begin{equation}
C_1C_2\le\tfrac{1+C_{d1}}{36}
\bigl(16C_{d2}k^{\beta}\eta^{2}+32KC_{d1}k^{2}\eta^{2}+16K\tauvar\eta^{3/2}\bigr),
\label{eq:formal-product}
\end{equation}
whose polynomial factor is $q^{2}$ rather than $q^{8}$; the declaration
\path{kernelProd_sqrt_}\hspace{0pt}\path{param_nolift} is this statement. The companion
shape \path{kernelProd_sqrt_}\hspace{0pt}\path{param} imposes $k^{2}B\le\tauvar$,
$k^{4}B\le\tauvar$ and collapses the two $\eta^{2}$ terms into a single $\tauvar\eta^{3/2}$ term
with the $\eta$-independent constant $(1+C_{d1})(16C_{d2}+32C_{d1}K+16K)/36$. The scale
$B_\sharp=\sqrt\eta$ carries no power of $q$.

\subsection{The asymptotic theorem and its two hypotheses}

The statements above are about the cost of a single admissible sequence and carry no side condition
beyond admissibility. The asymptotic corollary, by contrast, is phrased in the development through
Mathlib's real-valued \path{Filter.liminf} and therefore carries two representation-level
hypotheses: that an admissible sequence exists at every order
(\texttt{hne : $\forall$ N, (costs N $\varphi$).Nonempty}) and that the ratios $T_{\min}(N,\varphi)/N$
are eventually bounded above (\path{hbdd}). Both are artefacts of the definitions
(\path{sInf} of the empty set is $0$ in Mathlib, and \path{liminf} returns a junk value for
unbounded sequences), not assumptions of the analytic argument: the underlying contradiction is the
unconditional statement that no admissible sequence has $\tauvar\le\rho\,q$ for all sufficiently
large $N$. The development also proves the dichotomy form
(\path{c_ge_four_or_grows}), which drops \path{hbdd} and concludes that either the
coefficient is at least $4$ or the ratios diverge to $+\infty$.

\subsection{Scope of the verification}

Machine-checked: the elementary bound with its exact constant $2$, at every order and every
$0<\phi\le\pi$; the sharpened bound with its explicit threshold $N_0(a)$ for $\phi=\pi$; and the
combined form~\eqref{eq:formal-allN}, plus the lemmas they rest on. Not machine-checked:
the numerical material of Section~\ref{sec:numerics}. The low-order sequences found by local
search and their re-verified costs are upper bounds computed outside Lean; they enter no proof,
and no optimality or small-$N$ statement about them is formalized. The threshold $N_0(a)$ is
conservative by construction: $5010$ and $46000$ come from rounding constants up and from the
crude exponent combination~\eqref{eq:formal-bookkeeping}, not from the method, and the same
mechanism with numerically evaluated constants reaches $a\approx1.5$--$2$ at $q\approx10^{2}$.
The audit ran on the full development: the build is green (3739 jobs), a source scan finds no
\path{sorry}, \path{admit}, or \path{axiom} in the project files, and the eleven headline
theorems listed in \path{RobustZ/AuditFinal.lean} --- among them \path{endgame_bound},
\path{all_N_bound_end}, \path{elementary_bound_pi} and \path{collar_bound_concrete}
--- depend only on the three standard logical axioms \path{propext}, \path{Classical.choice}
and \path{Quot.sound}. Statements, build log and the raw \texttt{\#print axioms} output are
collected at \url{https://show.dytchem.cn/lean/}, and the repository is cited as~\cite{LeanRepo},
with sources at \url{https://github.com/Dytchem/robust-z-lower-bound-lean}.
